\documentclass[12pt]{article}

% === Packages ===
\usepackage[margin=1in]{geometry}
\usepackage{amsmath, amssymb}
\usepackage{setspace}
\usepackage{titlesec}
\usepackage{enumitem}
\usepackage{helvet}
\renewcommand{\familydefault}{\sfdefault} % modern sans-serif look

% === Formatting tweaks ===
\setstretch{1.15}
\emergencystretch=3em % avoid overfull lines in long references
\setlist[enumerate]{itemsep=0.4em, topsep=0.2em}
\titleformat{\section}{\Large\bfseries\sffamily}{\thesection.}{1em}{}
\titleformat{\subsection}{\large\bfseries\sffamily}{\thesubsection}{1em}{}

% === Title ===
\title{\vspace{-2em}\textbf{Applied Mathematics Oral Exam Questions and Model Answers}\\
\large Solid Mechanics Fundamentals}
\author{Prepared for Graduate Oral Examination}
\date{}

\begin{document}
\section*{Oral Exam Questions and Model Answers Based on the Solid Mechanics Tutorial}

%%%%%%%%%%%%%%%%%%%%%%%%%%%%%%%%%%%%%%%%%%%%%%%%%%%%%%%
\subsection*{1. What is stress, and how is it defined in 1D and 3D?}

\textbf{Answer.}
For a uniform bar under axial loading, the nominal normal stress is the applied force $F_{\text{applied}}$ divided by the original cross-sectional area $A$:
\[
\sigma = \frac{F_{\text{applied}}}{A}.
\]
(Here $F_{\text{applied}}$ is a scalar force; the symbol $\mathbf F_{\chi}$ is reserved for the deformation gradient in the questions below.)


In 3D, stress becomes a \textbf{second-order tensor} whose components 
$\sigma_{ij}$ represent the $i$-direction force acting on a surface with 
normal in the $j$-direction. Stress in 3D is written as
\[
\boldsymbol{\sigma} =
\begin{bmatrix}
\sigma_{11} & \sigma_{12} & \sigma_{13} \\
\sigma_{21} & \sigma_{22} & \sigma_{23} \\
\sigma_{31} & \sigma_{32} & \sigma_{33}
\end{bmatrix}.
\]
More generally, Cauchy's traction law is $\mathbf t(\mathbf n)=\boldsymbol\sigma\mathbf n$. The tensor is needed because the force per area depends on the orientation of the surface through its normal $\mathbf n$.


%%%%%%%%%%%%%%%%%%%%%%%%%%%%%%%%%%%%%%%%%%%%%%%%%%%%%%%
\subsection*{2. What is strain, and what does it measure?}

\textbf{Answer.}
The engineering strain of a uniformly stretched bar is
\[
\epsilon = \frac{\Delta L}{L_0}
\]


for small or moderate axial extension. In several dimensions there is no single strain tensor for every regime. The infinitesimal strain $\boldsymbol\varepsilon=\tfrac12(\nabla\mathbf u+\nabla\mathbf u^T)$ is appropriate for small deformation, whereas finite-deformation models commonly use the deformation gradient $\mathbf F_{\chi}=\nabla_X\boldsymbol\chi$ and objective measures such as $\mathbf E=\tfrac12(\mathbf F_{\chi}^T\mathbf F_{\chi}-\mathbf I)$. A proper strain measure records extension and shear while excluding rigid-body rotation.


%%%%%%%%%%%%%%%%%%%%%%%%%%%%%%%%%%%%%%%%%%%%%%%%%%%%%%%
\subsection*{3. Explain Hooke’s law.}

\textbf{Answer.}
In 1D, Hooke’s law is
\[
\sigma = E \epsilon,
\]
where $E$ is the Young’s modulus.

This is a linear, small-strain constitutive law. In three-dimensional isotropic linear elasticity, stress depends on the full strain tensor through two independent elastic constants; the scalar relation $\sigma=E\epsilon$ applies to a uniaxial test under the corresponding lateral conditions (lateral surfaces traction-free, so the bar is free to contract sideways by the Poisson effect). If lateral contraction is prevented instead, the axial stiffness is the larger constrained modulus $\lambda+2\mu$.


%%%%%%%%%%%%%%%%%%%%%%%%%%%%%%%%%%%%%%%%%%%%%%%%%%%%%%%
\subsection*{4. What types of stress are shown in the tutorial, and why do they matter in biomechanics?}

\textbf{Answer.}
In biomechanics:
\begin{itemize}
    \item \textbf{Tensile}: ligaments and tendons (stretching).
    \item \textbf{Compressive}: bone and cartilage (weight-bearing).
    \item \textbf{Shear}: tissue sliding, muscle layers, arterial walls.
\end{itemize}
These are loading modes rather than three different mathematical stress objects: a single stress tensor may contain tensile, compressive, and shear components. IBFE simulations often capture them simultaneously, especially in soft tissues subjected to combined loading.


%%%%%%%%%%%%%%%%%%%%%%%%%%%%%%%%%%%%%%%%%%%%%%%%%%%%%%%
\subsection*{5. Define traction and explain how it differs from pressure.}

\textbf{Answer.}
For a surface with unit normal $\mathbf n$, traction is the force per unit current area
\[
\mathbf t(\mathbf n)=\boldsymbol\sigma\mathbf n.
\]
It may have both normal and tangential components. Pressure is a scalar; the isotropic pressure part of the Cauchy stress is $-p\mathbf I$ and therefore contributes traction $-p\mathbf n$. Viscous or elastic stresses can add normal and shear traction. In IBFE, one must also distinguish current-area traction from the reference-area traction $\mathbf P\mathbf N$ associated with the first Piola--Kirchhoff stress.


%%%%%%%%%%%%%%%%%%%%%%%%%%%%%%%%%%%%%%%%%%%%%%%%%%%%%%%
\subsection*{6. What are the three main types of mechanical loading shown in the tutorial?}

\textbf{Answer.}
The tutorial uses three canonical loading modes:

\begin{itemize}
    \item \textbf{Axial}: produces uniform tensile or compressive stress.
    \item \textbf{Bending}: produces tension on one side and compression on the other.
    \item \textbf{Torsion}: produces shear stress.
\end{itemize}

In IBFE simulations, immersed structures (e.g. beams, muscles, tendons)
experience all of these depending on fluid forces and constraints.


%%%%%%%%%%%%%%%%%%%%%%%%%%%%%%%%%%%%%%%%%%%%%%%%%%%%%%%
\subsection*{7. Describe Young’s modulus, Poisson’s ratio, and yield strength.}

\textbf{Answer.}
The three quantities are defined as follows.
\begin{itemize}
    \item \textbf{Young's modulus $E$}: stiffness; slope of linear portion of the stress–strain curve.
    \item \textbf{Poisson's ratio $\nu$}: negative ratio of lateral to axial strain in a uniaxial small-strain test,
    \[
    \nu = -\epsilon_{\text{lateral}}/\epsilon_{\text{axial}}.
    \]
    \item \textbf{Yield strength $\sigma_y$}: stress at which irreversible deformation occurs.
\end{itemize}

In IBFE:
\begin{itemize}
    \item $E$ determines stiffness of structural elements and time-step limits (for explicit structural coupling the stable time step shrinks roughly like $h\sqrt{\rho/E}$, the time for an elastic wave to cross one element of size $h$).
    \item for an isotropic linear elastic solid, positive strain energy requires $-1<\nu<1/2$ (equivalently $\mu>0$ and $K=E/[3(1-2\nu)]>0$); most materials have $0<\nu<1/2$,
    \item many soft biological tissues have $\nu \approx 0.45$--$0.49$ (nearly incompressible); values of $\nu$ near $0.5$ indicate a nearly incompressible isotropic linear response; exact incompressibility is a constraint rather than simply setting $\nu=0.5$ in a displacement-only formulation,
    \item $\sigma_y$ is relevant only when the constitutive model includes yielding or permanent deformation.
\end{itemize}


%%%%%%%%%%%%%%%%%%%%%%%%%%%%%%%%%%%%%%%%%%%%%%%%%%%%%%%
\subsection*{8. Compare linear elastic, nonlinear elastic, viscoelastic, and viscoplastic materials.}

\textbf{Answer.}
The classes differ in whether the stress--strain response is linear, whether it is reversible, and whether it depends on time or loading rate:

\begin{itemize}
    \item \textbf{Linear elastic}: $\sigma \propto \epsilon$; material returns to original shape.
    \item \textbf{Nonlinear elastic}: stress--strain relation is nonlinear but reversible.
    \item \textbf{Hyperelastic}: a nonlinear elastic model derived from a strain-energy density $W$; it is especially convenient for finite deformation.
    \item \textbf{Viscoelastic}: time-dependent response (creep, stress relaxation).
    \item \textbf{Viscoplastic}: time-dependent permanent deformation.
\end{itemize}

IBFE must choose the appropriate constitutive law depending on the biological system:
soft tissues often require nonlinear or viscoelastic models rather than linear elasticity.


%%%%%%%%%%%%%%%%%%%%%%%%%%%%%%%%%%%%%%%%%%%%%%%%%%%%%%%
\subsection*{9. What boundary conditions are commonly used in solid mechanics?}

\textbf{Answer.}
At each boundary point and in each direction one prescribes either a displacement component or a traction component (exactly one of the pair). The common types of loading and constraint are:

\begin{enumerate}
    \item Prescribed displacement (an essential or Dirichlet condition)
    \item Prescribed pressure or traction (a natural or Neumann condition)
    \item Mixed displacement--traction conditions on different boundary portions
    \item Body forces, such as gravity, included as volume loads rather than boundary conditions
    \item Contact constraints and contact tractions between solids
    \item Thermal strains or other process-induced loads, with matching thermal and mechanical boundary data
\end{enumerate}

IBFE uses:
\begin{itemize}
    \item Lagrangian solid BCs to anchor or fix material points,
    \item Eulerian traction conditions for fluid–structure coupling,
    (more precisely: the fluid--structure interface is not a boundary of a body-fitted fluid mesh; the elastic force density, including the surface term from $\mathbf P\mathbf N$ on $\partial B_0$, is spread to the Eulerian grid as a body force, and traction balance at the interface emerges from the coupled equations rather than being imposed as a boundary condition),
    \item Contact or self-contact when structural meshes interact.
\end{itemize}


%%%%%%%%%%%%%%%%%%%%%%%%%%%%%%%%%%%%%%%%%%%%%%%%%%%%%%%
\subsection*{10. What is the importance of isotropy vs anisotropy in biological materials?}

\textbf{Answer.}
Isotropy means the constitutive response is unchanged by rotating the material; anisotropy means it depends on direction, typically because of aligned fibres or microstructure:

\begin{itemize}
    \item \textbf{Isotropic linear elastic}: material properties identical in all directions.
    \item \textbf{Anisotropic}: directional dependence (e.g., bone, tendon, myocardium).
\end{itemize}

In IBFE:
\begin{itemize}
    \item Isotropy simplifies constitutive laws.
    \item Anisotropy is required for fiber-reinforced tissues or materials whose stiffness differs along fiber directions.
\end{itemize}
Quantitatively, isotropic linear elasticity has 2 independent constants, transversely isotropic 5, orthotropic 9, and fully anisotropic up to 21; finite-strain fibre models add terms such as $W_f(I_4)$ with $I_4=\mathbf a_0\cdot\mathbf C\mathbf a_0$ for a fibre direction $\mathbf a_0$.



%%%%%%%%%%%%%%%%%%%%%%%%%%%%%%%%%%%%%%%%%%%%%%%%%%%%%%%
\subsection*{11. How do stress and strain generalize in 3D, and why do we need tensors?}

\textbf{Answer.}
For finite deformation, $\mathbf F_{\chi}$ maps reference line elements to current line elements. Objective strain tensors constructed from $\mathbf F_{\chi}$, such as $\mathbf E=\tfrac12(\mathbf F_{\chi}^T\mathbf F_{\chi}-\mathbf I)$, quantify stretch and shear without treating a rigid rotation as strain. Stress tensors map a surface normal to its traction. Thus:
\begin{itemize}
    \item Stress is a tensor describing forces acting on all possible surface orientations.
    \item Strain is a tensor describing how material line elements stretch and shear; rigid-body rotation is represented by the motion or deformation gradient, not by an objective strain measure.
\end{itemize}

Tensors allow:
\begin{itemize}
    \item compact representation of multidirectional relationships,
    \item coordinate-invariant properties,
    \item implementation in IBFE finite-element formulations where the deformation gradient and Cauchy stress appear naturally.
\end{itemize}

%%%%%%%%%%%%%%%%%%%%%%%%%%%%%%%%%%%%%%%%%%%%%%%%%%%%%%%
\subsection*{References for these answers}
\begin{itemize}
    \item A. F. Bower, \emph{Applied Mechanics of Solids}, CRC Press, 2010 (online at solidmechanics.org): Ch.~1 (boundary conditions), \S2.1--2.3, \S3.2 (linear elasticity).
    \item M. E. Gurtin, E. Fried, L. Anand, \emph{The Mechanics and Thermodynamics of Continua}, Cambridge University Press, 2010.
    \item G. A. Holzapfel, \emph{Nonlinear Solid Mechanics}, Wiley, 2000 (hyperelasticity, fibre-reinforced models).
    \item B. E. Griffith, X. Luo, Hybrid finite difference/finite element immersed boundary method, \emph{Int.\ J.\ Numer.\ Methods Biomed.\ Eng.} 33 (2017) e2888 (IBFE).
\end{itemize}

%%%%%%%%%%%%%%%%%%%%%%%%%%%%%%%%%%%%%%%%%%%%%%%%%%%%%%%
\end{document}
